% =====================================================================
%  Hoja de puntuacion -- una por persona evaluadora
%  Proyecto SIMPA -- Equipo AHMRV -- ISR-401 -- UTEQ -- 2026-2027 PPA
%  Instrumento del cuasi-experimento del Enfoque 1
%  Compilacion: pdflatex hoja_puntuacion.tex  (x2)
% =====================================================================
\documentclass[11pt,a4paper]{article}

\usepackage[utf8]{inputenc}
\usepackage[T1]{fontenc}
\usepackage[margin=1.8cm]{geometry}
\usepackage{array}
\usepackage{booktabs}
\usepackage{longtable}
\usepackage[table]{xcolor}
\usepackage{enumitem}
\usepackage{tcolorbox}
\usepackage{fancyhdr}
\usepackage{titlesec}
\usepackage{lastpage}

\definecolor{uteqverde}{HTML}{1F6F3C}
\definecolor{grisclaro}{HTML}{F2F2F2}

\newcolumntype{C}[1]{>{\centering\arraybackslash}p{#1}}
\newcolumntype{L}[1]{>{\raggedright\arraybackslash}p{#1}}

\titleformat{\section}{\normalfont\large\bfseries\color{uteqverde}}{\thesection.}{0.5em}{}
\titlespacing*{\section}{0pt}{10pt}{4pt}

\pagestyle{fancy}
\fancyhf{}
\renewcommand{\headrulewidth}{0.4pt}
\lhead{\footnotesize\color{uteqverde}Hoja de puntuación --- Evaluador/a \underline{\hspace{1.4cm}}}
\rhead{\footnotesize SIMPA --- Equipo AHMRV}
\cfoot{\footnotesize Página \thepage\ de \pageref{LastPage}}

\setlength{\parskip}{4pt}
\setlength{\parindent}{0pt}

\begin{document}

\begin{center}
{\footnotesize UNIVERSIDAD TÉCNICA ESTATAL DE QUEVEDO\\
Facultad de Ciencias de la Computación --- Carrera de Software}\\[8pt]
{\Large\bfseries\color{uteqverde} HOJA DE PUNTUACIÓN}\\[3pt]
{\footnotesize Evaluación de calidad de requisitos funcionales --- OSF \texttt{4z35d}}
\end{center}

\vspace{4pt}

\renewcommand{\arraystretch}{1.3}
\begin{center}
\begin{tabular}{L{3.2cm} L{5.5cm} L{2.6cm} L{4.2cm}}
\toprule
\rowcolor{grisclaro}
\textbf{Código evaluador/a} & \textbf{Fecha} & \textbf{Hora inicio} & \textbf{Hora fin} \\
\midrule
\rule{0pt}{4ex} & & & \\
\bottomrule
\end{tabular}
\end{center}

\begin{tcolorbox}[colback=grisclaro,colframe=uteqverde,boxrule=0.6pt,arc=2pt,left=4pt,right=4pt,top=3pt,bottom=3pt]
\footnotesize
\textbf{Escala:} 1 ausente \quad\textbar\quad 2 marginal \quad\textbar\quad
3 deficiencias apreciables \quad\textbar\quad 4 deficiencias menores
\quad\textbar\quad 5 plenamente presente.\\
\textbf{Dimensiones:} \textbf{COM} completitud \quad\textbar\quad
\textbf{AMB} ausencia de ambigüedad \quad\textbar\quad
\textbf{VER} verificabilidad \quad\textbar\quad
\textbf{COR} corrección respecto de la fuente \quad\textbar\quad
\textbf{CON} consistencia interna.\\
Puntúe las cinco dimensiones de cada requisito antes de pasar al siguiente. Si
duda entre dos niveles, elija el menor.
\end{tcolorbox}

\vspace{4pt}

\renewcommand{\arraystretch}{1.45}
\begin{longtable}{C{1.6cm} C{1.3cm} C{1.3cm} C{1.3cm} C{1.3cm} C{1.3cm} L{5.4cm}}
\toprule
\rowcolor{grisclaro}
\textbf{Req.} & \textbf{COM} & \textbf{AMB} & \textbf{VER} & \textbf{COR} &
\textbf{CON} & \textbf{Observaciones} \\
\midrule
\endfirsthead
\toprule
\rowcolor{grisclaro}
\textbf{Req.} & \textbf{COM} & \textbf{AMB} & \textbf{VER} & \textbf{COR} &
\textbf{CON} & \textbf{Observaciones} \\
\midrule
\endhead
\bottomrule
\endfoot
R-01 & & & & & & \\
R-02 & & & & & & \\
R-03 & & & & & & \\
R-04 & & & & & & \\
R-05 & & & & & & \\
R-06 & & & & & & \\
R-07 & & & & & & \\
R-08 & & & & & & \\
R-09 & & & & & & \\
R-10 & & & & & & \\
R-11 & & & & & & \\
R-12 & & & & & & \\
R-13 & & & & & & \\
R-14 & & & & & & \\
R-15 & & & & & & \\
R-16 & & & & & & \\
R-17 & & & & & & \\
R-18 & & & & & & \\
R-19 & & & & & & \\
R-20 & & & & & & \\
R-21 & & & & & & \\
R-22 & & & & & & \\
R-23 & & & & & & \\
R-24 & & & & & & \\
R-25 & & & & & & \\
R-26 & & & & & & \\
R-27 & & & & & & \\
R-28 & & & & & & \\
R-29 & & & & & & \\
R-30 & & & & & & \\
R-31 & & & & & & \\
R-32 & & & & & & \\
R-33 & & & & & & \\
R-34 & & & & & & \\
R-35 & & & & & & \\
R-36 & & & & & & \\
R-37 & & & & & & \\
R-38 & & & & & & \\
R-39 & & & & & & \\
R-40 & & & & & & \\
R-41 & & & & & & \\
R-42 & & & & & & \\
R-43 & & & & & & \\
R-44 & & & & & & \\
R-45 & & & & & & \\
R-46 & & & & & & \\
R-47 & & & & & & \\
R-48 & & & & & & \\
R-49 & & & & & & \\
R-50 & & & & & & \\
\end{longtable}

\newpage

\section*{Cierre de la sesión}

\renewcommand{\arraystretch}{1.4}
\begin{center}
\begin{tabular}{L{7.5cm} L{8.5cm}}
\toprule
\rowcolor{grisclaro}
\textbf{Pregunta} & \textbf{Respuesta} \\
\midrule
¿Consultó la rúbrica durante la sesión? & \\
¿Hubo algún requisito que no pudo puntuar? ¿Cuál? & \\
¿Comentó sus puntuaciones con otra persona evaluadora durante la sesión? & \\
Observaciones generales & \\[3ex]
\bottomrule
\end{tabular}
\end{center}

\vspace{10pt}

\begin{tcolorbox}[colback=white,colframe=uteqverde!45,boxrule=0.5pt,arc=2pt]
\footnotesize
La tercera pregunta no busca señalar a nadie. El protocolo exige puntuación
independiente, y si hubo intercambio de criterio debe declararse para
interpretar correctamente el grado de acuerdo entre evaluadores. Responder que
sí no invalida sus puntuaciones.
\end{tcolorbox}

\vspace{18pt}

\begin{center}
\begin{tabular}{C{7cm} C{7cm}}
\rule{6.5cm}{0.4pt} & \rule{6.5cm}{0.4pt} \\
{\footnotesize Firma de la persona evaluadora} &
{\footnotesize Firma de quien modera la sesión} \\
\end{tabular}
\end{center}

\vfill
{\footnotesize
Instrumento del protocolo registrado en OSF \texttt{4z35d}. Las puntuaciones se
digitalizan en \texttt{07\_Datos/datos\_crudos/evaluaciones\_experimento.csv} y
esta hoja se conserva como evidencia primaria.}

\end{document}
